% =====================================================================
%  CONJURA CONJECTURE STATEMENT
%  Tight security of the single-session lattice Chevallier-Mames
%  signature.
%  Requires conjura-conjecture.cls in the same directory.
% =====================================================================
\documentclass{conjura-conjecture}

\runninghead{SINGLE-SESSION LATTICE CHEVALLIER-MAMES}

\newcommand{\ZZ}{\mathbb{Z}}
\newcommand{\RR}{\mathbb{R}}
\newcommand{\NN}{\mathbb{N}}
\newcommand{\II}{\mathbb{I}}
\newcommand{\cR}{\mathcal{R}}
\newcommand{\cD}{\mathcal{D}}
\newcommand{\cC}{\mathcal{C}}
\newcommand{\cA}{\mathcal{A}}
\newcommand{\cB}{\mathcal{B}}
\newcommand{\cQ}{\mathcal{Q}}
\newcommand{\negl}{\mathrm{negl}}
\newcommand{\Rot}{\mathsf{Rot}}
\newcommand{\Adv}{\mathsf{Adv}}
\newcommand{\Time}{\mathsf{Time}}
\newcommand{\smlwe}{\mathsf{sMLWE}}
\newcommand{\SSone}{\mathsf{SS}^1}
\newcommand{\ssk}{\sigma_{\mathsf{sk}}}
\newcommand{\std}{\sigma_{\mathsf{td}}}
\newcommand{\esm}{\varepsilon_{\mathsf{sm}}}
\newcommand{\Hone}{\mathsf{H}_1}
\newcommand{\Htwo}{\mathsf{H}_2}

\begin{document}

\cjkicker{CONJURA \textperiodcentered{} OPEN PROBLEM}
\cjtitle{Tight Security of the Single-Session Lattice
  Chevallier-Mames Signature}
\cjsubtitle{One challenge per query instead of two, from search
  Module-LWE}
\cjstatus{Statement: AI-written from Rutchathon
  Chairattana-Apirom, Nico D\"ottling, Julian Loss, Stefano Tessaro,
  Benedikt Wagner, \emph{Tight Lattice-Based Signatures without
  Trapdoors from Search LWE}, Cryptology ePrint Archive, full version
  of an article to appear at CRYPTO 2026, 2026, not yet checked by a
  human. The two-session scheme $\mathsf{SS}_{\mathsf{R}}$ is settled
  --- the paper's Theorem~4.2 proves tight strong unforgeability for
  it from search Module-LWE --- while the single-session scheme
  $\SSone$ below is open: for it the paper sketches only a reduction
  that resamples challenges and therefore loses a factor about
  $1/\varepsilon$, and states that it does not know how to prove tight
  security.}
\cjcategory{\emph{Public-Key Cryptography \& Assumptions}.}

\vspace{0.4em}

\begin{informalconjecture}
There is a lattice signature scheme built by translating the
Chevallier-Mames CDH-based signature into the Module-LWE world: the
public key is a noisy linear sample, the signer derives a second matrix
from a hash of its own commitment, publishes a second noisy sample
under that matrix with the same secret, and proves that both samples
are close to the row span of their respective matrices with the same
secret. A security reduction wants to take a forgery and read the
secret out of it, and it can: from two different accepting challenges
for the same query it recovers the noise term and then the secret. The
trouble is the arithmetic of how often that happens. If a single
challenge is accepting with probability epsilon, then a pair of
accepting challenges shows up only with probability epsilon squared, so
the reduction has to keep resampling challenges, and the resulting
proof loses a polynomial factor. The paper's fix is to run two
independent copies of the protocol side by side under one public key
with two independent challenges; a short inclusion-exclusion argument
then shows that a pair of accepting challenges for one of the two
copies is at least as likely as a forgery, which makes the reduction
tight, at the cost of doubling signature size and signing time. The
question left open is whether the doubling is really needed: is the
single-copy scheme itself tightly secure under search Module-LWE, or is
the parallel repetition inherent to this style of proof?
\end{informalconjecture}

\section{The setting}\label{sec:setting}

A security reduction for a signature scheme is \emph{tight} when a
forger running in time $t$ with success probability $\varepsilon$
yields an algorithm for the underlying problem running in time
$t' \approx t$ with success probability $\varepsilon' \approx
\varepsilon$. Tightness matters concretely: a lossy reduction forces
larger parameters to absorb the loss. For lattices, the hash-and-sign
approach of Gentry, Peikert and Vaikuntanathan~\cite{GPV08} has a tight
reduction to SIS but requires trapdoor preimage sampling, while the
trapdoor-free Fiat--Shamir route of Lyubashevsky~\cite{Lyu12} is proved
by rewinding and is not amenable to a tight reduction to SIS; tight
proofs for Fiat--Shamir lattice signatures, including the one for
Dilithium~\cite{LDK22}, instead go through lossy identification and
therefore rely on \emph{decisional} LWE. Since known
search-to-decision reductions for LWE lose a factor polynomial in
$1/\varepsilon$, a proof from \emph{search} LWE is quantitatively
preferable, and obtaining one for an efficient Fiat--Shamir lattice
signature is the paper's stated goal.

The paper's starting point is the Chevallier-Mames
signature~\cite{Che05}, abstracted as a five-move identification
protocol in~\cite{KLP17}, which has a tight reduction to CDH: a
signature is a $\Sigma$-protocol transcript proving that $X = g^x$ and
$Z = h^x$ share a discrete logarithm, where $h$ is a hash of the
commitment and the message. The naive lattice translation is the scheme
$\SSone$ of Definition~\ref{def:ss1}: the public key is a Module-LWE
sample $\mathbf{t} = \mathbf{s}\mathbf{A} + \mathbf{e}$, the signer
derives $\mathbf{H} = \Hone(\mathbf{u},\mu)$ from its own commitment,
publishes $\bar{\mathbf{t}} = \mathbf{s}\mathbf{H} + \bar{\mathbf{e}}$,
and the verification equations assert simultaneous proximity of
$\mathbf{t}$ to $\mathbf{s}\mathbf{A}$ and of $\bar{\mathbf{t}}$ to
$\mathbf{s}\mathbf{H}$. Signing is simulatable without the secret key
by rejection sampling~\cite{Lyu12} together with a rerandomisation of
$(\mathbf{A},\mathbf{t})$ into
$(\mathbf{H},\bar{\mathbf{t}})$~\cite{GMPW20}; hardness of $\smlwe$ for
a uniform secret and a Gaussian error of this width follows from
worst-case module lattice problems~\cite{LS15}, provided
$\ssk > \sqrt{n}\cdot\omega(\sqrt{\log N})$.

The obstruction is on the extraction side. Unlike $\mathbf{e}$, the
term $\bar{\mathbf{e}}$ is adversarially chosen, and an adversary can
pick a non-short $\bar{\mathbf{e}}$ for which $c \cdot \bar{\mathbf{e}}$
is nonetheless short with inverse polynomial probability over the
challenge $c$: the proximity proof is not sound, so decoding
$\mathbf{s}$ from $\bar{\mathbf{t}}$ with a lattice trapdoor for
$\mathbf{H}$~\cite{MP12} fails. What survives is a shortcut: once $c$
is fixed, the response $\mathbf{z}$ is statistically determined, so a
reduction holding a trapdoor for $\mathbf{H}$ can decide from an
$\Htwo$ query alone whether the query can lead to a forgery, without
waiting for the forgery. From two accepting challenges $c \neq c'$ for
the same query one gets
\[
  (c-c')\mathbf{t} = (\mathbf{z}-\mathbf{z}')\mathbf{A}
    + \mathbf{f} - \mathbf{f}',
\]
whence $\mathbf{e} = (c-c')^{-1}(\mathbf{f}-\mathbf{f}')$ computed in
the number field, and then $\mathbf{s}$ by linear algebra. But if a
single challenge is accepting with probability $\varepsilon$, a
\emph{pair} of accepting challenges appears only with probability
$\varepsilon^2$, so the reduction must resample about $1/\varepsilon$
times, and the proof is not tight.

The paper's fix is a two-fold parallel repetition: two sessions share
one public key and one $\mathbf{H}$, and $\Htwo$ returns two
independent challenges $(c_0,c_1)$. Writing $\varepsilon_b$ for the
per-session acceptance probability, the inclusion-exclusion identity
\[
  \varepsilon_0^2 + \varepsilon_1^2 - \varepsilon_0^2\varepsilon_1^2
    \;\ge\; \varepsilon_0\varepsilon_1
\]
says that a pair of accepting challenges for \emph{one} of the two
sessions is at least as likely as an accepting pair $(c_0,c_1)$ for
both, i.e.\ at least as likely as the adversary's own success. This
``polarization'', combined with a balanced sampler that keeps the
reduction's \emph{worst-case} rather than expected running time under
control~\cite{JT20} and a monotone coupling argument lifting the
per-query guarantee to the whole game, yields a tight reduction from
$\smlwe$ to strong unforgeability of the two-session scheme, with
polynomial modulus. That proof factors through key-only unforgeability:
strong unforgeability under chosen-message attack follows from key-only
unforgeability together with $\smlwe$ (the paper's Section~4.3), and
the extraction argument above is what carries the key-only step. The
price of the repetition is exactly a factor two in signature size and
signing time. The question recorded here is whether that price is
necessary.

Progress to date. The paper gives, in its technical overview, a
non-tight reduction for the single-session scheme: the reduction
generates $\mathbf{H}$ with a lattice trapdoor, and because the
response is statistically determined once the challenge is fixed, it
can decide from each $\Htwo$ query whether the query is accepting; on
an accepting query it resamples challenges until a second accepting one
appears, which costs about $1/\varepsilon$ attempts, and then recovers
$\mathbf{e}$ and $\mathbf{s}$ from the two transcripts. It also shows
precisely what the two-session variant buys: for two independently
sampled challenges the probability of an accepting pair within one
session is at least the probability of an accepting pair across the two
sessions, which is the adversary's advantage. No such identity is
available with a single challenge per query, where the pair probability
is $\varepsilon^2$.

\section{Notation and parameters}\label{sec:notation}

Throughout, $\kappa$ is the security parameter, $N = N(\kappa)$ is a
power of two, $q = q(\kappa)$ is an odd prime, and
\[
  \cR := \ZZ[\mathsf{X}]/\langle \mathsf{X}^N + 1\rangle,
  \qquad \cR_q := \cR/q\cR.
\]
We write $\poly(\kappa)$ and $\negl(\kappa)$ for the classes of
polynomially bounded and of negligible functions of $\kappa$. All
vectors are row vectors. An element
$a = \sum_{i=0}^{N-1} a_i \mathsf{X}^i \in \cR$ is identified with its
coefficient vector $(a_0,\dots,a_{N-1}) \in \ZZ^N$, and a vector
$\mathbf{a} \in \cR^m$ with the concatenation of the coefficient
vectors of its entries, an element of $\ZZ^{mN}$; $\|\cdot\|$,
$\|\cdot\|_1$ and $\|\cdot\|_\infty$ denote the $\ell_2$, $\ell_1$ and
$\ell_\infty$ norms of that integer vector. For $a \in \cR$,
$\Rot(a) \in \ZZ^{N\times N}$ is the negacirculant matrix of
multiplication by $a$ on $\cR$, and for
$\mathbf{e} = (e_1,\dots,e_m) \in \cR^m$ we set
\[
  \Rot(\mathbf{e}) := (\Rot(e_1) \mid \cdots \mid \Rot(e_m))
    \in \ZZ^{N \times mN}.
\]

For $\sigma > 0$ and $\mathbf{v} \in \cR^m$,
$\cD_{\cR^m,\mathbf{v},\sigma}$ denotes the discrete Gaussian on
$\cR^m$ (through the above embedding) with standard deviation $\sigma$
centred at $\mathbf{v}$, i.e.\ the distribution with probability mass
proportional to $\exp(-\pi\|\mathbf{x}-\mathbf{v}\|^2/\sigma^2)$; we
write $\cD_{\cR^m,\sigma} := \cD_{\cR^m,\mathbf{0},\sigma}$ and denote
by $\cD_{\cR^m,\mathbf{v},\sigma}(\mathbf{x})$ the probability it
assigns to $\mathbf{x}$. For a symmetric positive definite
$\boldsymbol\Sigma \in \RR^{mN\times mN}$,
$\cD_{\cR^m,\sqrt{\boldsymbol\Sigma}}$ is the discrete Gaussian on
$\cR^m$ centred at $\mathbf{0}$ with covariance $\boldsymbol\Sigma$,
i.e.\ with probability mass proportional to
$\exp(-\pi\,\mathbf{x}\boldsymbol\Sigma^{-1}\mathbf{x}^\top)$. We write
$\II_k$ for the $k\times k$ identity matrix and $\otimes$ for the
Kronecker product.

For $\gamma_c \in \NN$, the challenge set is
\[
  \cC := \{\, c \in \cR \;:\;
    \|c\|_\infty = 1 \ \wedge\ \|c\|_1 = \gamma_c \,\},
\]
the set of ring elements with coefficients in $\{-1,0,1\}$ of which
exactly $\gamma_c$ are nonzero. We write $x \sample S$ for sampling $x$
uniformly from a finite set $S$ or according to a distribution $S$.

The parameters of the statement below are as follows.

\begin{itemize}[leftmargin=1.2em]
  \item $\kappa$: security parameter.
  \item $N$: degree of the cyclotomic ring, a power of two.
  \item $q$: odd prime modulus of the ring $\cR_q$.
  \item $n, m$: row and column dimensions of the public matrix
  $\mathbf{A} \in \cR_q^{n\times m}$; also the dimensions of the search
  Module-LWE instance.
  \item $m'$: number of columns of the hash-derived matrix
  $\mathbf{H} \in \cR_q^{n\times m'}$.
  \item $\ssk$: standard deviation of the discrete Gaussian
  secret-key error $\mathbf{e}$.
  \item $\sigma_r, \bar\sigma_r$: standard deviations of the commitment
  noise $\mathbf{d}$ and $\bar{\mathbf{d}}$, and the target widths of
  the rejection sampling.
  \item $\sigma_x, \sigma_y$: parameters of the covariance
  $\bar{\boldsymbol\Sigma} =
   \sigma_x^2\Rot(\mathbf{e})\Rot(\mathbf{e})^\top + \sigma_y^2\II_N$
  of the second error $\bar{\mathbf{e}}$.
  \item $\beta_{\mathbf{z}}, \bar\beta_{\mathbf{z}}$: verification norm
  bounds for the two proximity checks.
  \item $\gamma_c$: number of nonzero coefficients of a challenge,
  i.e.\ the $\ell_1$ norm of every $c \in \cC$.
  \item $M$: rejection sampling parameter,
  $M = \exp(t/\alpha + 2/\alpha^2)$; for $\SSone$ the per-iteration
  acceptance probability is $1/M^2$, so the expected number of signing
  iterations is about $M^2$ (the paper's two-session scheme has
  $1/M^4$).
  \item $\alpha, t$: the two rejection-sampling slack parameters
  entering $M$, both $\omega(\sqrt{\log(mN)})$.
  \item $\ell$: maximum number of rejection sampling iterations before
  signing gives up.
  \item $\esm$: the smoothing-parameter slack appearing in the bounds
  on $\ssk$ and $\eta$ (written $\varepsilon$ in the paper's Table~1;
  renamed here to keep $\varepsilon$ for a forger's success
  probability).
  \item $\eta$: the Gaussian width below which the paper's
  rerandomisation and trapdoor arguments apply.
  \item $b, \std$: the gadget base and the trapdoor Gaussian width of
  the trapdoor sampler used by the reduction; they constrain $m'$ and
  $q$ only.
  \item $\cQ_{\mathsf{S}}, \cQ_1, \cQ_2$: numbers of signing, $\Hone$
  and $\Htwo$ queries made by the forger.
  \item $C_0, p$: the multiplicative advantage loss and the per-query
  time overhead polynomial that make the reduction tight.
\end{itemize}

\section{The conjecture}\label{sec:conjectures}

\begin{definition}[Search Module-LWE]
\label{def:smlwe}
Let $n,m,N,q$ be as above and $\sigma > 0$. For an algorithm $\cB$
define
\[
  \Adv^{\mathrm{smlwe}}_{n,m,N,q,\sigma}(\cB,\kappa)
    := \Pr[\, \mathbf{s} = \mathbf{s}' \,],
\]
where the probability is over
$\mathbf{A} \sample \cR_q^{n\times m}$,
$\mathbf{s} \sample \cR_q^{n}$,
$\mathbf{e} \sample \cD_{\cR^m,\sigma}$ and
$\mathbf{s}' \sample \cB(\mathbf{A},\, \mathbf{s}\mathbf{A}+\mathbf{e})$,
with $\mathbf{A}$ and $\mathbf{s}$ uniform and all arithmetic over
$\cR_q$. The assumption $\smlwe_{n,m,N,q,\sigma}$ states that this
advantage is negligible for every probabilistic polynomial-time $\cB$.
\end{definition}

\begin{definition}[The single-session scheme $\SSone$]
\label{def:ss1}
Let $n,m,m',\ell,\gamma_c \in \NN$, $M \ge 1$, and
$\ssk,\sigma_r,\bar\sigma_r,\sigma_x,\sigma_y,\beta_{\mathbf{z}},
\bar\beta_{\mathbf{z}} > 0$, and let
$\Hone : \{0,1\}^* \to \cR_q^{n\times m'}$ and
$\Htwo : \{0,1\}^* \to \cC$ be hash functions. The signature scheme
\begin{multline*}
  \SSone = \SSone[n,m,m',N,q,\ell,M,\ssk,\sigma_r,\\
  \bar\sigma_r,\sigma_x,\sigma_y,\beta_{\mathbf{z}},
  \bar\beta_{\mathbf{z}},\gamma_c,\Hone,\Htwo]
\end{multline*}
consists of the following algorithms, with all arithmetic over $\cR_q$.

\smallskip\noindent$\mathsf{KeyGen}(1^\kappa)$: sample
$\mathbf{A} \sample \cR_q^{n\times m}$, $\mathbf{s} \sample \cR_q^{n}$
and $\mathbf{e} \sample \cD_{\cR^m,\ssk}$, set
$\mathbf{t} \leftarrow \mathbf{s}\mathbf{A} + \mathbf{e}$, and return
$\mathsf{pk} := (\mathbf{A},\mathbf{t})$ and
$\mathsf{sk} := (\mathsf{pk},\mathbf{s},\mathbf{e})$.

\smallskip\noindent$\mathsf{Sign}(\mathsf{sk},\mu)$ for
$\mu \in \{0,1\}^*$: set
\[
  \bar{\boldsymbol\Sigma} := \sigma_x^2\,\Rot(\mathbf{e})
    \Rot(\mathbf{e})^\top + \sigma_y^2\,\II_N \in \RR^{N\times N}.
\]
Repeat the following at most $\ell$ times; if no signature is output,
return $\bot$.
\begin{enumerate}[leftmargin=1.4em]
\item Sample $\mathbf{r} \sample \cR_q^{n}$,
$\mathbf{d} \sample \cD_{\cR^m,\sigma_r}$ and
$\bar{\mathbf{d}} \sample \cD_{\cR^{m'},\bar\sigma_r}$, and set
$\mathbf{u} \leftarrow \mathbf{r}\mathbf{A} + \mathbf{d} \in \cR_q^{m}$.
\item Set
$\mathbf{H} \leftarrow \Hone(\mathbf{u},\mu) \in \cR_q^{n\times m'}$.
\item Sample
$\bar{\mathbf{e}} \sample
  \cD_{\cR^{m'},\sqrt{\II_{m'}\otimes\bar{\boldsymbol\Sigma}}}$
and set
$\bar{\mathbf{t}} \leftarrow \mathbf{s}\mathbf{H} + \bar{\mathbf{e}}
 \in \cR_q^{m'}$ and
$\bar{\mathbf{u}} \leftarrow \mathbf{r}\mathbf{H} + \bar{\mathbf{d}}
 \in \cR_q^{m'}$.
\item Set
$c \leftarrow \Htwo(\mathbf{u},\bar{\mathbf{u}},\mathbf{H},
 \bar{\mathbf{t}},\mu) \in \cC$.
\item Set $\mathbf{z} \leftarrow \mathbf{r} + c\,\mathbf{s}$,
$\mathbf{f} \leftarrow \mathbf{d} + c\,\mathbf{e}$ and
$\bar{\mathbf{f}} \leftarrow \bar{\mathbf{d}} + c\,\bar{\mathbf{e}}$
(over $\cR$).
\item With probability $\rho\cdot\bar\rho$ return
$\sigma := (\bar{\mathbf{t}},\mathbf{u},\bar{\mathbf{u}},\mathbf{z})$;
otherwise go to the next iteration. Here
\[
  \rho := \min\Bigl\{1,
    \tfrac{\cD_{\cR^{m},\sigma_r}(\mathbf{f})}
          {M\,\cD_{\cR^{m},c\mathbf{e},\sigma_r}(\mathbf{f})}\Bigr\}
\]
and
\[
  \bar\rho := \min\Bigl\{1,
    \tfrac{\cD_{\cR^{m'},\bar\sigma_r}(\bar{\mathbf{f}})}
          {M\,\cD_{\cR^{m'},c\bar{\mathbf{e}},\bar\sigma_r}
            (\bar{\mathbf{f}})}\Bigr\}.
\]
\end{enumerate}

\smallskip\noindent$\mathsf{Ver}(\mathsf{pk},\mu,\sigma)$: parse
\[
  \sigma = (\bar{\mathbf{t}},\mathbf{u},\bar{\mathbf{u}},\mathbf{z})
    \in \cR_q^{m'}\times\cR_q^{m}\times\cR_q^{m'}\times\cR_q^{n},
\]
set $\mathbf{H} \leftarrow \Hone(\mathbf{u},\mu)$ and
$c \leftarrow \Htwo(\mathbf{u},\bar{\mathbf{u}},\mathbf{H},
\bar{\mathbf{t}},\mu)$, and return $1$ if and only if both
\begin{gather*}
  \|\mathbf{u} + c\cdot\mathbf{t} - \mathbf{z}\mathbf{A}\|
    \le \beta_{\mathbf{z}}, \\
  \|\bar{\mathbf{u}} + c\cdot\bar{\mathbf{t}} - \mathbf{z}\mathbf{H}\|
    \le \bar\beta_{\mathbf{z}},
\end{gather*}
where the two vectors are taken with representatives in
$\{-\lfloor q/2\rfloor,\dots,\lfloor q/2\rfloor\}$ coefficientwise.
\end{definition}

\begin{definition}[Correctness and $\mathrm{SUF\text{-}CMA}$ security]
\label{def:sufcma}
A signature scheme $\mathsf{SS}$ has correctness error at most
$\varepsilon$ if for every $\mu \in \{0,1\}^*$,
$\Pr[\mathsf{SS}.\mathsf{Ver}(\mathsf{pk},\mu,\sigma) = 0] \le
\varepsilon$, over
$(\mathsf{sk},\mathsf{pk}) \sample \mathsf{SS}.\mathsf{KeyGen}(1^\kappa)$
and $\sigma \sample \mathsf{SS}.\mathsf{Sign}(\mathsf{sk},\mu)$ (an
output of $\bot$ counts as a failure). For an adversary $\cA$ with
access to a signing oracle $\mathsf{S}$ that on input $\mu$ returns
$\sigma \sample \mathsf{SS}.\mathsf{Sign}(\mathsf{sk},\mu)$ and records
the pair $(\mu,\sigma)$ in a set $\mathsf{Sigs}$, define
\[
  \Adv^{\mathrm{sufcma}}_{\mathsf{SS}}(\cA,\kappa) :=
  \Pr\bigl[\, (\mu^*,\sigma^*) \notin \mathsf{Sigs}
    \ \wedge\ \mathsf{Win} \,\bigr],
\]
where $\mathsf{Win}$ is the event
$\mathsf{SS}.\mathsf{Ver}(\mathsf{pk},\mu^*,\sigma^*) = 1$ and the
probability is over
$(\mathsf{sk},\mathsf{pk}) \sample
 \mathsf{SS}.\mathsf{KeyGen}(1^\kappa)$ and
$(\mu^*,\sigma^*) \sample \cA^{\mathsf{S}}(\mathsf{pk})$. In the random
oracle model $\cA$ additionally has oracle access to $\Hone$ and
$\Htwo$, drawn uniformly from the corresponding function spaces, and so
do $\mathsf{Sign}$ and $\mathsf{Ver}$.
\end{definition}

\begin{definition}[Admissible parameters]
\label{def:params}
An \emph{admissible parameter family} is a tuple of functions
\begin{multline*}
  \Pi = (n,m,m',N,q,\ell,\gamma_c,b,\\
  M,\ssk,\sigma_r,\bar\sigma_r,\sigma_x,\sigma_y,\\
  \beta_{\mathbf{z}},\bar\beta_{\mathbf{z}},\std,\alpha,t,\esm)
\end{multline*}
of $\kappa$, the first eight taking values in $\NN$ and the rest in
$\RR_{>0}$, with $\esm \in (0,1)$, all computable in time
$\poly(\kappa)$, such that $N$ is a power of two, $q$ is an odd prime,
$b \ge 2$, $M \ge 1$, all of $n,m,m',\ell,\gamma_c,\log q$ are bounded
by $\poly(\kappa)$, and such that for every $\kappa$ the following hold,
where
\[
  \eta := \tfrac{8}{\sqrt{\pi}}\, q^{n/m}
    \sqrt{N \ln\bigl(2mN(1+1/\esm)\bigr)}.
\]
\begin{itemize}[leftmargin=1.2em]
  \item $m > n$ and $m' \ge m + n\log_b q$;
  \item $\ssk \ge \sqrt{n \log N}$ and
  $\ssk \ge \sqrt{\log(2mN(1+1/\esm))/\pi}$;
  \item $\sigma_x = \sqrt{2}\,\eta$ and
  $\sigma_y = \sigma_x\,\ssk\,N\sqrt{m}$;
  \item $\sigma_r = \alpha\,\gamma_c\,\ssk\sqrt{mN}$ and
  $\bar\sigma_r = \alpha\,\gamma_c\,\sigma_y\sqrt{2m'N}$;
  \item $\beta_{\mathbf{z}} = \sigma_r\sqrt{mN}$ and
  $\bar\beta_{\mathbf{z}} = \bar\sigma_r\sqrt{m'N}$;
  \item $M = \exp(t/\alpha + 2/\alpha^2)$, with
  $\alpha = \omega(\sqrt{\log(mN)})$ and
  $t = \omega(\sqrt{\log(mN)})$;
  \item $|\cC| \ge 2^{2\kappa}$;
  \item $\std \ge \eta$ and
  $q \ge 4\,b\,\std\,\bar\beta_{\mathbf{z}}\sqrt{m'mN}$.
\end{itemize}
These are the paper's Table~1 constraints with the parallel repetition
removed.
\end{definition}

\begin{conjecture}[tight security of $\SSone$]\label{conj:main}
Let $\Pi$ be an admissible parameter family as in
Definition~\ref{def:params}, and let $\Hone$ and $\Htwo$ be modelled as
random oracles. Then the scheme $\SSone = \SSone[\Pi,\Hone,\Htwo]$ of
Definition~\ref{def:ss1} satisfies both of the following.

\smallskip\noindent\emph{(i) (Correctness.)} The correctness error of
$\SSone$ is $\negl(\kappa)$.

\smallskip\noindent\emph{(ii) (Tight security.)} There exist a constant
$C_0 \ge 1$ and a polynomial $p$ such that for every adversary $\cA$
against $\mathrm{SUF\text{-}CMA}$ security of $\SSone$ in the random
oracle model, running in time $\Time(\cA)$ and making at most
$\cQ_{\mathsf{S}}$ signing queries and at most $\cQ_1$ and $\cQ_2$
queries to $\Hone$ and $\Htwo$, there is an adversary $\cB$ against
$\smlwe_{n,m,N,q,\ssk}$ with
\[
  \Time(\cB) \;\le\; \Time(\cA)
    + \bigl(\cQ_{\mathsf{S}}\ell + \cQ_1 + \cQ_2\bigr)\cdot p(\kappa)
\]
and
\begin{align*}
  \Adv^{\mathrm{sufcma}}_{\SSone}(\cA,\kappa) \;\le\;
    & C_0 \cdot
      \Adv^{\mathrm{smlwe}}_{n,m,N,q,\ssk}(\cB,\kappa) \\
    & {}+ \negl(\kappa).
\end{align*}
In particular, $C_0$ and $p$ do not depend on $\cA$, and the loss is
independent of $\cQ_{\mathsf{S}},\cQ_1,\cQ_2$ and of
$\Adv^{\mathrm{sufcma}}_{\SSone}(\cA,\kappa)$.
\end{conjecture}

\begin{remark}[the constant $C_0$]\label{rem:cnought}
The paper's own bound for the two-session scheme
$\mathsf{SS}_{\mathsf{R}}$ has the two-adversary form
$\Adv^{\mathrm{smlwe}}(\cB) + \Adv^{\mathrm{smlwe}}(\cB') +
\delta_{\negl}$, from its Lemmas~4.4 and~4.5, so $C_0 = 2$ is the
faithful constant to aim at; the statement above leaves $C_0$ an
unspecified constant because that is what tightness requires.
\end{remark}

\begin{remark}[readings fixed by this write-up]\label{rem:readings}
The sentence the paper leaves this open with refers to ``this
protocol'', meaning the overview-level Attempt~1 on page~4, which omits
the rejection sampling and all parameter choices. It is instantiated
here as the one-session specialisation of the fully specified scheme in
Figure~2, which is the faithful reading (the overview says the missing
details are a routine modification), but a reader should check that
specialisation against Figure~2 on page~19. Two further deliberate
choices. The conjecture does not require a polynomial modulus, which
the paper does achieve for the two-session scheme, since that is not
part of what the paper says is open; and the formulation ``there exists
a reduction with constant loss'' is the standard theorem shape in this
literature and is not formalised by the paper itself, so anyone
attacking the negative direction will need to fix a meta-reduction
model. Finally, it has not been checked whether follow-up work has
since settled this; the paper is dated 2026 and only its own reference
list was available.
\end{remark}

\begin{remark}[why the statement is clean]\label{rem:clean}
The object is the paper's own published scheme with one bit of the
construction changed: run one session instead of two. Both the scheme
and the assumption are fully specified in the paper, so the statement
can be written from scratch in a page, and the target is a single
well-defined implication rather than a research programme. The
obstruction is named and quantitative: with one challenge per query, an
accepting pair has probability $\varepsilon^2$ rather than
$\varepsilon$, and the paper's inclusion-exclusion identity, which
repairs this, has no single-session analogue.
\end{remark}

\begin{remark}[why it is interesting]\label{rem:interest}
The question isolates a clean structural point: is two-fold parallel
repetition inherent to tight lattice Fiat--Shamir proofs from search
assumptions, or an artifact of this particular argument? A positive
answer halves the signature size and signing time of the first
efficient lattice Fiat--Shamir signature with a tight reduction to a
search assumption, on a scheme whose parameters are already large. A
negative answer, in the form of a meta-reduction ruling out tight
reductions for the single-session scheme, would be the first evidence
that the repetition is necessary and would sit naturally beside the
known impossibility results for tight reductions for Schnorr-type
signatures. Either way the answer says something about the polarization
technique the paper introduces, which is its main conceptual
contribution and is claimed to be of independent interest.
\end{remark}

\section{Bibliography}

\begin{conjurabibliography}{99}

\bibitem[Che05]{Che05}
B. Chevallier-Mames.
\emph{An efficient CDH-based signature scheme with a tight security
reduction}.
CRYPTO 2005, LNCS 3621, pages 511--526, Springer, 2005.

\bibitem[GMPW20]{GMPW20}
N. Genise, D. Micciancio, C. Peikert, M. Walter.
\emph{Improved discrete gaussian and subgaussian analysis for lattice
cryptography}.
PKC 2020, Part I, LNCS 12110, pages 623--651, Springer, 2020.

\bibitem[GPV08]{GPV08}
C. Gentry, C. Peikert, V. Vaikuntanathan.
\emph{Trapdoors for hard lattices and new cryptographic constructions}.
40th ACM STOC, pages 197--206, ACM Press, 2008.

\bibitem[JT20]{JT20}
J. Jaeger, S. Tessaro.
\emph{Expected-time cryptography: Generic techniques and applications
to concrete soundness}.
TCC 2020, Part III, LNCS 12552, pages 414--443, Springer, 2020.

\bibitem[KLP17]{KLP17}
E. Kiltz, J. Loss, J. Pan.
\emph{Tightly-secure signatures from five-move identification
protocols}.
ASIACRYPT 2017, Part III, LNCS 10626, pages 68--94, Springer, 2017.

\bibitem[LDK22]{LDK22}
V. Lyubashevsky, L. Ducas, E. Kiltz, T. Lepoint, P. Schwabe,
G. Seiler, D. Stehl\'e, S. Bai.
\emph{CRYSTALS-DILITHIUM}.
Technical report, National Institute of Standards and Technology, 2022.

\bibitem[LS15]{LS15}
A. Langlois, D. Stehl\'e.
\emph{Worst-case to average-case reductions for module lattices}.
Designs, Codes and Cryptography 75(3):565--599, 2015.

\bibitem[Lyu12]{Lyu12}
V. Lyubashevsky.
\emph{Lattice signatures without trapdoors}.
EUROCRYPT 2012, LNCS 7237, pages 738--755, Springer, 2012.

\bibitem[MP12]{MP12}
D. Micciancio, C. Peikert.
\emph{Trapdoors for lattices: Simpler, tighter, faster, smaller}.
EUROCRYPT 2012, LNCS 7237, pages 700--718, Springer, 2012.

\end{conjurabibliography}

\end{document}
